In the $k$-step-off pipeline, the update frequency $k$ controls the trade-off between summary staleness and prefilling overhead: smaller $k$ provides fresher summaries but requires more frequent uncached prefills, while larger $k$ amortizes this cost at the expense of increased lag. We ablate this hyperparameter under CPU offload with the \texttt{GLM-4.7} agent backbone, reporting results in Table~\ref{tab:k_sensitivity}.

\begin{table}[t]
\centering
\setlength{\tabcolsep}{3.0pt}
% \renewcommand{\arraystretch}{1.3}
\caption{Sensitivity to update frequency $k$ on SWE-bench Verified with \texttt{GLM-4.7} (CPU offload enabled).}
\label{tab:k_sensitivity}
\begin{tabular}{l c c c}
\toprule
Frequency $k$ & \%Resolved & Latency (s) & Speedup \\
\midrule
Full-Context & 69.0 & 5869.4 & 1.00$\times$ \\
\midrule
$k=16$ & 60.2 & 3576.8 & 1.64$\times$ \\
$k=8$ & 62.4 & 2836.0 & 2.07$\times$ \\
$k=4$ (default) & 62.7 & 2821.4 & 2.08$\times$ \\
$k=2$ & 64.2 & 2841.4 & 2.07$\times$ \\
$k=1$ & 57.2 & 3843.5 & 1.53$\times$ \\
\bottomrule
\end{tabular}
\vspace{-2mm}
\end{table}

Performance is stable across $k \in \{2, 4, 8\}$, with $k=2$ slightly outperforming the default $k=4$ in resolve rate. The two extremes exhibit notably degraded behavior: $k=1$ incurs excessive uncached prefilling overhead (as illustrated in Figure~\ref{fig:gantt}), while $k=16$ saturates the memory of the agent server and introduces scheduling overhead. Importantly, all {\model} variants with $k \in \{2, 4, 8\}$ achieve approximately $2\times$ wall-clock speedup over the full-context baseline while retaining the majority of its resolve rate, confirming that the method is not overly sensitive to this choice within a reasonable operating range.
